\section*{Chapitre \ref{ch:b1:s-block} --- L'hydrogène et le bloc s}

\begin{solution}{exo:b1:s-block:1}
\ce{KH} et \ce{CaH2}~: salins~; \ce{SiH4} et \ce{H2S}~: covalents~;
\ce{PdH_{0.6}}~: métallique. \ce{CaH2 + 2H2O -> Ca(OH)2 + 2H2}.
\end{solution}

\begin{solution}{exo:b1:s-block:2}
\ce{2Li + 2H2O -> 2LiOH + H2}, et de même avec Na et K~;
\ce{Mg + 2HCl -> MgCl2 + H2}.
\end{solution}

\begin{solution}{exo:b1:s-block:3}
Basiques~: \ce{Na2O}, \ce{MgO} (\ce{MgO + 2HCl -> MgCl2 + H2O}).
Acides~:
\ce{SiO2}, \ce{SO3}, \ce{CO2} (\ce{SO3 + H2O -> H2SO4}). Amphotères~:
\ce{Al2O3}, \ce{ZnO} (\ce{ZnO + 2OH- + H2O -> Zn(OH)4^2-}).
\end{solution}

\begin{solution}{exo:b1:s-block:4}
À pH 14, le couple de l'eau est \ce{2H2O + 2e- -> H2 + 2OH-}, $E = -0.83$~V~;
$\log K = 2(-0.83 + 2.71)/0.059 = 64$.
\end{solution}

\begin{solution}{exo:b1:s-block:5}
Dans la flamme, les chocs portent un électron sur un niveau excité~; en
revenant à l'\omterm{def:b1:quantum-numbers:energy-level}{état fondamental}, l'atome émet un photon d'énergie égale à
l'écart entre les niveaux, propre à l'élément. $E = hc/\lambda = 6.626 \times 10^{-34} \times
3.00 \times 10^8/(589 \times 10^{-9}) = \qty{3.375e-19}{J} = \qty{2.11}{eV}$.
% ledger: const:h, const:c, const:eV
\end{solution}

\begin{solution}{exo:b1:s-block:6}
$10^6/106.0 = \qty{9.43e3}{mol}$ de carbonate de sodium~:
\qty{9.43e3}{mol} de calcaire, soit \qty{0.944}{t}, et $2 \times 9.43 \times 10^3/0.90 =
\qty{2.10e4}{mol}$ de sel, soit \qty{1.23}{t}.
\end{solution}

\begin{solution}{exo:b1:s-block:7}
$[\ce{Mg^2+}] = 1.3/24.3 = \qty{0.053}{mol/L}$~; $\mathrm{pH} = 14.00 -
\frac12(11.25 + \log 0.053) = 9.0$. Les \qty{53.5}{mol} de magnésium de
\qty{1.0}{m^3} exigent \qty{53.5}{mol} de \ce{Ca(OH)2}~: \qty{4.0}{kg}.
\end{solution}

\begin{solution}{exo:b1:s-block:8}
La constante d'équilibre (thermodynamique) est énorme dans les deux
cas~; la vitesse (cinétique) n'est pas fixée par $E^\circ$. Le lithium
fond à plus haute température, reste solide et se recouvre d'une couche
de produit moins soluble, et perd son électron moins facilement que le
potassium (énergie d'ionisation plus élevée)~: il réagit
régulièrement, le potassium violemment.
\end{solution}

\begin{solution}{exo:b1:s-block:9}
$10^6/100.1 = \qty{9.99e3}{mol}$~: CaO $9.99 \times 10^3 \times 56.1 =
\qty{0.560}{t}$~; \ce{CO2} \qty{9.99e3}{mol}, $V = nRT/p = 9.99 \times 10^3
\times 8.314 \times 298.15/10^5 = \qty{248}{m^3}$.
\end{solution}

\begin{solution}{exo:b1:s-block:10}
$2.9 \times 10^8/6.9 = \qty{4.2e7}{kmol}$ de Li (masse en kg), $\times 73.8/2 =
\qty{1.55e9}{kg}$~: environ 1.6 million de tonnes de carbonate de
lithium.
$2.9 \times 10^8/8 = \num{3.6e7}$ batteries.
\end{solution}

\begin{solution}{exo:b1:s-block:11}
1.31, 0.84 et 0.71 pour cent à 20, 80 et \qty{100}{\celsius}~: la
\omterm{def:b1:precipitation:solubility}{solubilité} diminue. À chaud, moins de carbonate reste en solution~: on
en récupère davantage.
\end{solution}

\begin{solution}{exo:b1:s-block:12}
\ce{NaHCO3} est le moins soluble des sels que ces ions peuvent former, et
la solution en est saturée la première. L'ammoniac rend la solution
assez basique pour transformer \ce{CO2} en \ce{HCO3-} tout en restant
récupérable (sous forme de \ce{NH4Cl}, régénéré par la chaux)~;
l'hydroxyde de sodium serait consommé et coûterait plus cher que le
produit.
\end{solution}

\begin{solution}{pb:b1:s-block:1}
\textbf{1.} Li \qty{6.0}{g/L}, Mg \qty{2.4}{g/L}.
\textbf{2.} Quatre mètres cubes de saumure en donnent un~: trois mètres
cubes d'eau s'évaporent.
\textbf{3.} Le soleil et le climat sec fournissent l'énergie
gratuitement.
\textbf{4.} Le chlorure de sodium, de loin le sel le plus abondant, est
le premier à saturer.
\textbf{5.} Li $6000/6.9 = \qty{870}{mol}$~; Mg $2400/24.3 = \qty{98.8}{mol}$.
\textbf{6.} Le magnésium précipiterait avec le carbonate et
contaminerait le produit.
\textbf{7.} \ce{Mg^2+ + Ca(OH)2 -> Mg(OH)2 + Ca^2+}.
\textbf{8.} $98.8 \times 74.1 = \qty{7.32}{kg}$.
\textbf{9.} Il reste $[\ce{Mg^2+}] = \num{9.88e-5}$ mol/L~: $\mathrm{pH} = 14.00 -
\frac12(11.25 - 4.01) = 10.38$.
\textbf{10.} L'hydroxyde de lithium est soluble~: aucun solide ne se
forme à ce pH.
\textbf{11.} \ce{Ca^2+ + CO3^2- -> CaCO3}.
\textbf{12.} $98.8 \times 106.0 = \qty{10.5}{kg}$.
\textbf{13.} \ce{2Li+ + CO3^2- -> Li2CO3}.
\textbf{14.} $435 \times 106.0 = \qty{46.1}{kg}$.
\textbf{15.} $435 \times 73.8 = \qty{32.1}{kg}$.
\textbf{16.} Le carbonate de lithium est moins soluble à chaud.
\textbf{17.} $0.0084 \times 1000 = \qty{8.4}{kg}$.
\textbf{18.} Avec de petits volumes d'eau chaude, dans laquelle il est le
moins soluble.
\textbf{19.} $32.1 - 8.4 = \qty{23.7}{kg}$.
\textbf{20.} $23.7/32.1 = \qty{74}{\percent}$.
\textbf{21.} Elle contient encore du lithium~: on la renvoie dans les
bassins.
\textbf{22.} $2.9 \times 10^8/6.0 = \num{4.8e7}$ mètres cubes.
\textbf{23.} \textbf{\qty{24}{kg}} (23.7) de carbonate de lithium par
mètre cube.
\end{solution}
